\subsection{L2: the strengthened adic axiom does not imply condition 7}

\begin{theorem}[Refuted: L2, in every positive characteristic]
\label{A:L2-counterexample}
Let $A_0=\F_p[[t]]$ and $S=\Spec A_0$. On affine $S$-schemes define
\[
 H(R)=\{x\in R:x^p=0\text{ and }x\in tRx\}
     =\{x\in R:x^p=0,\ \exists a\in R\ (x=tax)\}.       \tag{A.2.3}
\]
This is an fpqc sheaf, commutes with filtered ring colimits, satisfies
the finite universal Artinian representability of condition 5,
valuative uniqueness, and condition 8 on locally finite-type tests.
It satisfies the strengthened condition 4 for every complete
noetherian adic ring, with no restriction on its ideal of definition.
Nevertheless it fails condition 7, with an actual modular pair
$A=B=A_0$ and $\xi=0$. Thus L2 is false even for a subgroup functor
of the finite flat group scheme $\alpha_{p,S}$.
\end{theorem}
\begin{proof}
The displayed definition is functorial. It gives a subgroup of
$\alpha_p(R)$: if $(1-ta)x=(1-tb)y=0$, then
\[
 (1-t(a+b-tab))(x+y)=(1-ta)(1-tb)(x+y)=0.
\]
Negatives and scalar multiples satisfy the same condition. The
following descent proof also defines its unique extension from
affines to an fpqc sheaf on schemes.

\emph{Conditions 1 and 2, including arbitrary fpqc descent.}
Let $R\to R'$ be faithfully flat, and let an element of $H(R')$
satisfy the descent equality. Descent for the affine scheme
$\alpha_p$ gives a unique $x\in R$ with $x^p=0$.
Put $M=Rx\subseteq R$. Flatness identifies
$M\otimes_R R'$ with the submodule $R'x$ of $R'$ and identifies
$(M/tM)\otimes_R R'$ with $R'x/tR'x$. Membership of the descended
element in $H(R')$ says precisely that the image of $[x]\in M/tM$
is zero. Faithful flatness implies $[x]=0$, or $x\in tRx$.
This proves the equalizer over every ring. Zariski locality follows
either by this argument on finite affine covers, or directly by
local membership in the quasi-coherent ideal $(tx)$.

\emph{Condition 3.}
The equations $x^p=0$ and $x=tax$ involve just the two elements
$x,a$ and finitely many ring operations. A section and its witness
therefore descend through a filtered ring colimit to a finite
stage. Equality of two descended sections is equality of their
underlying ring elements, which is also detected at a finite
stage. This proves the required bijection on filtered colimits.

\emph{The strengthened condition 4.}
Let $R$ be noetherian, complete and separated for an arbitrary
ideal $I$. A compatible family $x_n\in H(R/I^n)$ gives a unique
$x\in R$ with $x^p=0$. Choose a lift $a_n\in R$ of a witness
for $x_n$; no compatibility of these witnesses is assumed. Then
\[
 x-ta_nx\in I^n,\qquad
 x\in (tx)+I^n\quad\hbox{for every }n.
\]
Completeness implies $I\subseteq\operatorname{Jac}(R)$, by the
convergent geometric series for $1-i$, $i\in I$. Krull's
intersection theorem on the finite $R$-module $R/(tx)$,
in the Jacobson-radical form \cite[Lemma 10.51.5(2)]{StacksKrullJac}, gives
\[
 \bigcap_{n\geq1}\bigl((tx)+I^n\bigr)=(tx).
\]
Thus $x\in(tx)$, which is exactly the missing membership
condition. Injectivity follows from $I$-adic separatedness of
$R$. We have proved
\[
 H(R)\xrightarrow{\sim}\varprojlim_n H(R/I^n).
\]
This also gives the comparison with $H(\widehat R)$ for every
noetherian adic pair, by applying the result to its completion.

\emph{Condition 5, universally over each Artinian base.}
If $C$ is Artinian local, $t$ is either a unit or nilpotent.
In the first case $H_C=\alpha_{p,C}$ on \emph{all} $C$-algebras:
take $a=t^{-1}$. In the second case $1-ta$ is always a unit,
so $H_C$ is the terminal functor on all $C$-algebras. The
representing schemes are respectively
$\Spec C[X]/(X^p)$ and $\Spec C$. They are finite and local.
This proves the universal finite Artinian requirement, not just
a tangent-space assertion. Products treat nonlocal Artinian
bases if they are included in the convention.

\emph{Condition 6.}
A valuation ring is reduced, so its only element with $p$th
power zero is zero. Hence $H$ has just one point on valuation
rings and their fraction fields, and satisfies valuative
uniqueness.

\emph{Condition 8.}
Let $S'$ be irreducible and noetherian over $S$. If $t$ lies in
the generic prime, it is locally nilpotent, and $H_{S'}$ is
terminal. Otherwise the dense open $D(t)$ contains the generic
point and $H|_{D(t)}=\alpha_p$. In both cases a separated finite
representation is available on an open containing the given
generic point. For a test $Z$ locally of finite type, actual
local modularity at that point spreads to a nonempty open by
the local-ring argument in Proposition~\ref{A:8-literal}.
This is exactly the finite-type version of condition 8 in the
question.

\emph{Failure of the actual condition 7.}
Take $A=B=A_0$, and take the zero section $\xi$.
If $R$ is any local $S$-algebra centered at the closed point,
then $t\in\mathfrak m_R$, so $1-ta$ is a unit for every $a$.
Consequently $H(R)=\{0\}$. There is exactly one local
$S$-morphism $\Spec R\to\Spec A_0$, and its induced map onto
the eligible sections is therefore a bijection. The residue
field $\F_p$ is the minimal field of definition of the unique
closed-fibre point. Tests centered at the generic point have
no residue-field specialization of that closed-fibre point,
and hence no eligible section in this local universal property.
Thus $(\Spec A_0,0)$ is \emph{actually} locally modular in the
sense of cote 131, scan 35. Choosing the original base to be
$\Spec A_0$ also removes any possible ambiguity about whether
modularity is taken before or after the base change in 7.

At $\mathfrak p=(0)$ put $K=\F_p((t))$. Since $t$ is
invertible, $H_K=\alpha_{p,K}$. On the local test
$K[\epsilon]/(\epsilon^2)$ the zero section has the nonzero
eligible deformation $\epsilon$, since $\epsilon^p=0$.
The unique local $K$-map to $\Spec K$ induces only zero.
Thus $\xi_{\mathfrak p}$ is not locally modular. All the
hypotheses of 7 are present and its conclusion fails.
\end{proof}

The counterexample makes the correction precise: completeness for
all noetherian adic rings does not supply generization of actual
local modularity, even with the other axioms and a finite flat
ambient group. A separate generization hypothesis remains necessary.
Our positive replacement is the complete representability theorem
with the explicitly stronger pro-modular condition in
Theorem~\ref{A:pro7-representability}; the present example violates
both versions of 7. It also explains why the closed-diagonal
argument proposed for L2 cannot by itself establish that propagation.
